\section{Benchmarking PCPOP}
\label{sec:benchmark}
We report the performance of \texttt{PCPOP} on examples considered before. We use the size of the principal moment matrix, number of scalar variables and cpu time as measure of performance on a problem at specific level of relaxation.

We start with the example \ref{ex:pcpop_hardy} as it can be generalized to any number of parties to show genuine multipartite entanglement. We consider   

\begin{table}[H]
	\centering
	\caption{Benchmark summary for \texttt{PCPOP} relaxations. SDP size refers to the principal moment matrix dimension; CPU time measured in seconds.}
	\label{tab:pcpop-bench}
	\begin{tabular}{@{}ccrrr@{}}
		\# of parties & level & SDP size & \# scalar vars& CPU time (s) \\
		\midrule
		% Fill rows below with your data
		% parties  order  size  vars  time(s)
		2 & 2 & 13 & 31 & 0.127 \\
		2 & 3 & 25 & 61 & 0.612 \\
		2 & 4 & 41 & 101 & 1.917 \\
		3 & 2 & 25 & 93 & 0.516 \\
        3 & 3 & 63 & 250 & 5.189 \\
        3 & 4 & 129 & 529 & 45.915 \\
        4 & 2 & 41 & 229 & 1.646 \\
        4 & 3 & 129 & 833 & 44.58 \\
        4 & 4 & 321 & 2241 & 643.277 \\ 
		\bottomrule
	\end{tabular}
\end{table}
Now, we compare the performance of \texttt{PCPOP} with \texttt{Ncpol2sdpa} on the examples $\ref{ex:3comms},\ref{ex:4comms},\ref{ex:5comms}$. We compare the size of the principal moment matrix and the number of scalar variables for both packages at the same relaxation order.
\begin{table}[H]
	\centering
	\caption{Comparison of \texttt{PCPOP} and \texttt{Ncpol2sdpa} by relaxation order. SDP size refers to the principal moment matrix dimension.}
	\label{tab:compare-pcpop-ncpol}
	\begin{tabular}{@{}cccrrr@{}}
		\toprule
		Example & level & Gr\"obner size & Package & SDP size & \# scalar vars \\
		\midrule
		\multirow{2}{*}{\ref{ex:3comms}} & \multirow{2}{*}{2} & \multirow{2}{*}{6} & PCPOP      & 9 & 23 \\
		                           &                    &                      & Ncpol2sdpa & 9 & 25 \\
		\addlinespace
		\multirow{2}{*}{\ref{ex:3comms}} & \multirow{2}{*}{3} & \multirow{2}{*}{8} & PCPOP      & 17 & 58 \\
		                           &                    &                      & Ncpol2sdpa & 18 & 73 \\
		\addlinespace
        \multirow{2}{*}{\ref{ex:3comms}} & \multirow{2}{*}{4} & \multirow{2}{*}{10} & PCPOP      & 30 & 142 \\
		                           &                    &                      & Ncpol2sdpa & 35 & 223 \\
		\addlinespace

		\multirow{2}{*}{\ref{ex:4comms}} & \multirow{2}{*}{2} & \multirow{2}{*}{10} & PCPOP      & 15 & 60 \\
		                           &                    &                      & Ncpol2sdpa & 15 & 64 \\
		\addlinespace
		\multirow{2}{*}{\ref{ex:4comms}} & \multirow{2}{*}{3} & \multirow{2}{*}{14} & PCPOP      & 37 & 219 \\
		                           &                    &                      & Ncpol2sdpa & 40 & 284 \\
		\addlinespace
        \multirow{2}{*}{\ref{ex:4comms}} & \multirow{2}{*}{4} & \multirow{2}{*}{18} & PCPOP      & 82 & 723 \\
		                           &                    &                      & Ncpol2sdpa & 96 & 1339 \\
		\addlinespace
		\multirow{2}{*}{\ref{ex:5comms}} & \multirow{2}{*}{2} & \multirow{2}{*}{13} & PCPOP      & 24 & 176 \\
		                           &                    &                      & Ncpol2sdpa & 24 & 184 \\
		\addlinespace
		\multirow{2}{*}{\ref{ex:5comms}} & \multirow{2}{*}{3} & \multirow{2}{*}{19} & PCPOP      & 85 & 1686 \\
		                           &                    &                      & Ncpol2sdpa & 89 & 1934 \\
		\addlinespace
		\multirow{2}{*}{\ref{ex:5comms}} & \multirow{2}{*}{4} & \multirow{2}{*}{25} & PCPOP      & 287 & 17226 \\
		                           &                    &                      & Ncpol2sdpa & 315 & 22126 \\
		\bottomrule
	\end{tabular}
\end{table}
On top of improvements in the size of the principal moment matrix and number of scalar variables, \texttt{PCPOP} also provides tighter bounds as compared to \texttt{Ncpol2sdpa}. This is obvious as \texttt{PCPOP} imposes more constraints on th emoment matrix.

